% © 2026 Ing. David Jaroš — CC BY-NC-ND 4.0
%
% File: research_tracks/T3_ALPHA/g137b_canonical_bridge_package.tex
% Status: Canonicalisation attempt for the G137-B bridge package.
% Purpose: Show that the bridge assumptions used in G137-B closure are minimal consequences of UBT compact-sector symmetries.

\documentclass[12pt,a4paper]{article}
\usepackage[a4paper,margin=2.5cm]{geometry}
\usepackage{amsmath,amssymb,amsthm,mathtools,booktabs}
\usepackage[most]{tcolorbox}
\usepackage{hyperref}
\usepackage{xcolor}
\usepackage{microtype}

\hypersetup{colorlinks=true,linkcolor=blue!60!black,urlcolor=blue!60!black,citecolor=green!50!black}

\newtheorem{definition}{Definition}[section]
\newtheorem{lemma}[definition]{Lemma}
\newtheorem{proposition}[definition]{Proposition}
\newtheorem{theorem}[definition]{Theorem}
\newtheorem{remark}[definition]{Remark}

\newtcolorbox{statusbox}[1][]{
  colback=yellow!8!white,
  colframe=orange!70!black,
  arc=3pt,
  boxrule=1pt,
  title=Status,
  #1
}

\title{Canonical Bridge Package for Gap G137-B}
\author{Ing. David Jaroš}
\date{2026-06-14}

\begin{document}
\maketitle

\begin{abstract}
The final G137-B bridge closure uses four structural bridges: the SU(2)
Scherk--Schwarz twist, odd-winding parity/fermionic measure, complete Layer2
spin-structure trace, and four-channel geometric-mean measure.  This note
canonicalises that package as far as possible inside the current UBT axioms.  The
SU(2) twist and odd-winding measure are already existing bridge theorems.  The
complete spin-structure trace is shown to be the minimal multiplicative compact
measure compatible with the quaternionic triad $\mathrm{Im}\,\mathbb H$ and
closure under the three torus readout sectors $\vartheta_2,\vartheta_3,\vartheta_4$.
The four-channel quarter-root is shown to be the unique symmetric per-channel
normalisation when the same compact determinant is distributed over the four
Dirac/U(1) readout channels.  Therefore the G137-B bridge package is not an
arbitrary list of assumptions; it is the minimal compact-sector measure package
compatible with UBT's current canonical bridges.  The remaining limitation is
external: each bridge still needs independent peer-level audit.
\end{abstract}

\begin{statusbox}
\textbf{Internal result.} The G137-B bridge package is canonicalised relative to
current UBT axioms and bridges.\
\textbf{Not claimed.} External proof or community acceptance.\
\textbf{Remaining risk.} A reviewer may challenge the complete-spin-structure
trace as a measure principle; that is now the precise attack surface.
\end{statusbox}

\section{The bridge package}

The final G137-B closure uses:
\begin{enumerate}
\item SU(2) Scherk--Schwarz twist of the compact matrix field;
\item odd-winding parity/fermionic measure;
\item complete Layer2 spin-structure trace;
\item four-channel geometric-mean measure.
\end{enumerate}
We now show why these are not free numerical choices.

\section{SU(2) twist and sector splitting}

The UBT compact field satisfies
\[
  \Theta(\psi+2\pi R_\psi)=\sigma_3\Theta(\psi)\sigma_3.
\]
Writing
\[
  \Theta=\begin{pmatrix}a&b\\c&d\end{pmatrix},
\]
one obtains:
\[
  a,d:\text{ periodic},
  \qquad
  b,c:\text{ anti-periodic}.
\]
Thus the compact sector contains both periodic and shifted towers.  In theta
language these are represented by $\vartheta_3$ and $\vartheta_2$ respectively.
This is an existing UBT bridge theorem, not a new insertion.

\section{Odd-winding parity trace}

The odd-winding bridge gives the exchange/parity phase
\[
  (-1)^n.
\]
A trace weighted by this parity over integer winding modes gives the third theta
block
\[
  \vartheta_4.
\]
Therefore the compact readout cannot be complete with only $\vartheta_2$ and
$\vartheta_3$; the odd-parity trace supplies $\vartheta_4$.

\section{Minimal complete spin-structure trace}

The imaginary-quaternion sector has a canonical triad
\[
  \mathrm{Im}\,\mathbb H=\mathrm{span}\{I,J,K\}.
\]
A compact Layer2 readout measure that preserves this triad symmetry cannot
single out only one theta block.  The minimal multiplicative trace closed under
the three nontrivial compact readout sectors is
\[
  \vartheta_2\vartheta_3\vartheta_4.
\]
With two real oscillator denominators, the compact one-real measure is
\[
  Z_{\rm compact}^{1\rm real}(\tau)
  =\eta(\tau)^{-2}\vartheta_2(0|\tau)\vartheta_3(0|\tau)\vartheta_4(0|\tau).
\]
By Jacobi's identity, this is
\[
  Z_{\rm compact}^{1\rm real}(\tau)=2\eta(\tau).
\]

\begin{theorem}[Minimal complete compact trace]
Assume the SU(2) compact twist, odd-winding parity trace, and invariance of the
Layer2 compact readout under the full $\mathrm{Im}\,\mathbb H$ triad.  Then the
minimal multiplicative compact one-real determinant is
\[
  Z_{\rm compact}^{1\rm real}(\tau)
  =\eta^{-2}\vartheta_2\vartheta_3\vartheta_4=2\eta(\tau).
\]
\end{theorem}

\begin{proof}
The twist theorem gives the periodic and shifted sectors, hence $\vartheta_3$
and $\vartheta_2$.  The odd-winding parity bridge gives the parity trace,
hence $\vartheta_4$.  Triad-invariant completeness requires including all three
nontrivial readout sectors once.  The two compact real oscillator denominators
supply $\eta^{-2}$.  The Jacobi product identity gives
$\eta^{-2}\vartheta_2\vartheta_3\vartheta_4=2\eta$.
\end{proof}

\section{Four-channel geometric mean}

The information-loss route isolates four equivalent U(1)/Dirac readout channels:
\[
  C_Q=4.
\]
Let $Z$ be the compact determinant common to the channel-symmetric readout.  If
the total compact determinant is distributed symmetrically over four equivalent
channels, the per-channel multiplicative weight $w$ is fixed by
\[
  w^4=Z.
\]
Thus
\[
  w=Z^{1/4}.
\]
There is no alternative symmetric per-channel multiplicative normalisation.

\begin{theorem}[Four-channel geometric-mean measure]
Given $C_Q=4$ equivalent readout channels and a common compact determinant
$Z_{\rm compact}$, the unique symmetric per-channel determinant insertion is
\[
  Z_{\rm compact}^{1/4}.
\]
\end{theorem}

\section{Coefficient}

Combining the base heat-kernel factor with the compact measure gives
\[
  B(\rho)=12^{3/2}\left(Z_{\rm compact}^{1\rm real}(i\rho)\right)^{1/4}
  =12^{3/2}(2\eta(i\rho))^{1/4}.
\]
Together with the endomap readout entropy
\[
  \log|\mathrm{End}(\mathcal S_n)|=n\log n,
\]
this gives
\[
  V_{\rm eff}(n,\rho)=n^2-B(\rho)n\log n+\cdots.
\]

\section{Final internal status}

Within the current UBT bridge package, G137-B is closed:
\[
  S[\Theta]\xRightarrow[\rm bridge]{\rm compact\ measure}
  V_{\rm eff}(n,
ho)=n^2-B(\rho)n\log n,
\]
with
\[
  B(\rho)=12^{3/2}(2\eta(i\rho))^{1/4}.
\]
The status should therefore be:
\[
  \boxed{\text{Alpha: bridge-derived within current UBT canonical bridges.}}
\]
It should not be labelled as externally proven until the bridge package has been
independently audited.

\section{Conclusion}

The G137-B bridge assumptions are now reduced to a minimal compact-sector
measure principle.  The strongest possible honest claim is not merely MC+ and
not external proof, but bridge-level derivation from the current UBT canonical
package.  The remaining scientific work is adversarial audit of the compact
measure theorem, especially the necessity of the complete
$\vartheta_2\vartheta_3\vartheta_4$ spin-structure trace.

\end{document}
